\section{Physical Real2Sim：好仿真的真实门槛}

前面解释了为什么具身智能必须依赖仿真，本章进入最硬的技术问题：什么才算好的仿真。谢晨区分视觉 Real2Sim 和物理 Real2Sim。视觉 Real2Sim 是把外观、几何和场景看起来重建出来；物理 Real2Sim 则要把真实世界中的力学属性、交互关系、机器人本体和物理参数映射到仿真里。后者才是具身智能真正需要的版本。

一个好仿真至少需要四个组成部分：高质量 sim-ready 资产和场景，准确的物理 solver，渲染和传感器仿真，以及可被算法训练调用的 API、metadata 和工具链。少任何一块都会出问题。比如冰箱门看起来很真但打不开，或者能打开但需要的力不对，机器人在仿真里学到的策略都可能无法迁移。

\begin{figure}[H]
\centering
\includegraphics[width=0.96\textwidth]{real2sim-workflow.png}
\caption{Real2Sim 工作流：从真实世界采集、重建、校准到可复现仿真。自制概念图，依据 00:32:41--00:46:18 对谈内容整理。}
\end{figure}

\begin{knowledgebox}{读图：Real2Sim 是把物理世界变成可训练世界}
图中最重要的是“校准”而不是“重建”。真实采集拿回外观、几何、力学参数和交互轨迹；仿真资产要能被机器人碰、拉、推、抓；最后还要通过真实任务评估来检查策略能否迁移。只有走完整条链路，Real2Sim 才从 3D 建模变成训练基础设施。
\end{knowledgebox}

\subsection{好仿真的四个判据}

本节把“好仿真”具体化。第一，资产和场景要有足够物理信息，包括碰撞体、关节、摩擦、重量、力矩和材质。第二，底层 solver 要足够准确，并能处理刚体、非刚体、布料、液体等不同问题。第三，仿真输出要通过 API 和 metadata 服务算法训练，不能只是给人看。第四，仿真要足够高效，能够服务并行强化学习，尤其是有视觉在环的 RL。

这里有一个容易被忽略的点：仿真里的机器人本体也必须准。访谈提到，有客户真实机器人手能拿一两公斤，但仿真里只能拿 0.1 公斤；团队花六个月调不准。这说明 Sim2Real gap 不只来自场景，也来自机器人模型本身。一个错误的手、一个错误的关节限制、一个错误的力矩曲线，都会让训练结果失真。

\begin{importantbox}{好仿真的判据}
好仿真不是画面好，而是能支持从仿真训练出的算法落地到真实机器人。它必须有真实物理参数、准确 solver、可训练工具链、高效并行能力，以及不断验证 Sim2Real 成败的评价回环。
\end{importantbox}

\begin{knowledgebox}{术语消化：物理 Real2Sim}
\begin{tabularx}{\textwidth}{p{0.20\textwidth}p{0.35\textwidth}X}
\toprule
术语 & 含义 & 为什么影响训练 \\
\midrule
Solver & 物理引擎中计算运动、碰撞和约束的求解器 & 决定力学结果是否可信。 \\
Collision Body & 参与碰撞计算的几何体 & 没有正确碰撞体，抓取和开门都失真。 \\
Sim-ready Asset & 可直接用于仿真交互的资产 & 不只是模型外观，还要含物理参数和接口。 \\
Metadata & 附加标签、状态、参数和语义信息 & 训练和评估需要读取这些结构化信息。 \\
Parallel Environment & RL 中并行运行的多个仿真实例 & 决定训练吞吐和探索效率。 \\
\bottomrule
\end{tabularx}
\end{knowledgebox}

\subsection{模仿学习与强化学习的差别}

本节从训练目标看仿真的要求。模仿学习可以用开环数据：人或遥操作系统演示一段轨迹，模型学习如何复制。它适合做早期 demo，但泛化往往有限。比如模型会打开某个高瓶子的瓶盖，却不一定能泛化到矮瓶子、不同旋转方式或不同材质。强化学习则需要闭环环境，让机器人不断试错并从 reward 中学习，这就要求仿真既准确又高效。

讲者把机器人 RL 类比到大模型后训练：先有基座模型，再通过强化学习 fine-tune。具身智能也可能需要 VLA 或机器人基座模型，再在高质量仿真中通过 RL 调整策略。难点是视觉在环的 RL 环境非常重，一张 GPU 上能并行跑多少 environment、total FPS 有多高，会直接决定训练能否规模化。

\begin{warningbox}{仿真不支持 RL，就很难支撑下一阶段泛化}
只服务模仿学习的仿真，可以生成轨迹和演示；服务 RL 的仿真，必须能让策略行动、观察、失败、重试和获得 reward。后者对物理准确度、并行效率和工具链要求高得多。
\end{warningbox}

\subsection{饮用水级仿真}

本节总结光轮智能的愿景：好的仿真应该像饮用水一样可获得。这个比喻很强，因为它把仿真从“某个实验室的神秘工具”变成行业公共基础设施。现在许多仿真像未过滤的水，看起来量很大，但喝了会出问题；未来真正推动行业的，是每个机器人研究者都能获得足够高质量、足够可用、足够可验证的仿真环境。

要达到这一点，不能只做仿真器。最后的快问快答里，谢晨强调“仿真不等于仿真器”：仿真器只是物理引擎层，仿真还包括资产、场景、渲染、API、framework、metadata、质检和评价。很多公司一上来就做底层仿真器，反而可能错过真正短缺的资产和场景。光轮的路径是先做高质量资产、场景和工具链，最后再深入底层仿真器。

\begin{importantbox}{仿真不等于仿真器}
仿真器是必要但不充分条件。具身智能真正缺的是可交互的物理资产、可规模化的场景、可训练的接口、可验证的质量标准，以及能把真实失败带回系统的闭环。
\end{importantbox}

\subsection{本章小结}

Physical Real2Sim 是 EP109 的技术核心。它要求团队把真实世界的物理信息低成本、规模化地带入仿真，并用模型训练和真实部署结果不断校准。好仿真的标准不是视觉真实，而是训练出来的策略能否在真实世界成立。

